\appendixchapter{Permutations and Determinants over a Field}
\label{app:determinants}
\chaptermark{Permutations and Determinants}
\section*{Introduction}
This appendix supplements the self-contained determinant construction in
Chapter~10 with a slower treatment of permutations and an arbitrary-field
viewpoint. The inversion-based sign proof in the main text needs no appendix;
the parity argument below is an alternative conceptual proof. A field $F$ has addition, subtraction, multiplication, and
division by nonzero elements obeying the familiar arithmetic laws, with
$0\ne1$. No background in abstract algebra is assumed. Unless otherwise
stated, vector spaces here are finite dimensional over $F$. Replace real
scalars by elements of $F$ in the vector-space axioms. Taking $F=\R$
recovers the main text; neither this appendix nor the next is a prerequisite
for Chapters~11--13.
The arbitrary-field determinant will also supply the coordinate detectors
for exterior bases in Appendix~\ref{app:tensors}; its value lies in the
field even when no notion of positive orientation is available.
The permutation exposition is influenced by the author's algebra
notes~\cite{chen-algebra}, specialized here to the elementary facts needed
for determinants.

\section{What Survives over an Arbitrary Field?}
Span, independence, exchange, bases, dimension, linear maps, matrices, dual
spaces, and multilinear maps use only field arithmetic. The Chapter~10
proofs therefore work over $F$. Determinants, tensor products, and exterior
powers also belong to this algebraic theory, as the constructions below show.

A compatible total order, positivity, absolute value or norm, metric,
topology, convergence, and orientation via determinant sign are additional
structures, not consequences canonically supplied by field axioms.
Order, positivity, norms, topology, and convergence must be supplied
separately. For instance, a symmetric bilinear form over a field need not
be positive definite; the field axioms do not even specify a meaning for
``positive.'' The algebraic constructions here need none of these additions.

Integer coefficients encode repeated addition: $m$ in $F$ means
$m\cdot1_F$ for $m\geq0$, and negative integers mean additive inverses.
The arithmetic laws survive, but distinct integers can have the same value.
For a first example $\mathbb F_2=\{0,1\}$ has $1+1=0$, $1\cdot1=1$,
and $1^{-1}=1$. These tables obey the field laws. Here $3=1$ and $-2=0$.

\begin{example}[Finite fields and characteristic two]
For prime $p$, the residues $\mathbb F_p$ with arithmetic modulo $p$
form a field, and $\mathbb F_p^n$ has the usual standard basis.
For completeness, elementary integer arithmetic supplies the inverse:
if $p\nmid a$, Bezout's identity gives integers $u,v$ with $ua+vp=1$,
so the residue of $u$ inverts that of $a$. Bezout's identity is used here
as elementary arithmetic background; no quotient-ring theory is needed.
The \textbf{characteristic} of a field is the least positive integer $p$
such that $p\cdot1=0$, or zero if there is no such integer.
A positive characteristic is prime: if $p=ab$ with $1<a,b<p$,
then $(a\cdot1)(b\cdot1)=0$ contradicts invertibility of nonzero factors.
An ordered field has $1>0$; repeated addition would give $p\cdot1>0$,
contradicting $p\cdot1=0$. Thus positive-characteristic fields cannot
be ordered fields, and ordinary Euclidean positivity is unavailable.
Over $\mathbb F_2^2$, the symmetric bilinear expression
$B(x,y)=x_1y_1+x_2y_2$ has
$B((1,1),(1,1))=1+1=0$ although $(1,1)\ne0$.
\[
\boxed{\text{alternating means vanishing on repeated inputs}.}
\]
This is the field-independent definition. Expansion of
two equal inputs $u+v$ still proves the interchange identity $A(u,v)=-A(v,u)$
(with any other slots fixed). In characteristic two, $-1=1$, so an
alternating bilinear form is also symmetric. Symmetry alone does not force
zero diagonal: $B(x,y)=xy$ on $\mathbb F_2$ has $B(1,1)=1$.
The same warning applies to exterior products: $v\wedge v=0$ is fundamental;
the consequence $v\wedge w=-w\wedge v$ becomes $v\wedge w=w\wedge v$.
\end{example}

\section{Cycles, Transpositions, and Parity}
A permutation of $\{1,\ldots,n\}$ is a bijection of this set to itself.
Their set is denoted $S_n$. Composition is composition of bijections:
$(\sigma\tau)(j)=\sigma(\tau(j))$. We read the right-hand permutation
first. The identity fixes every entry, and $\sigma^{-1}$ reverses all
arrows of $\sigma$. No formal group theory is needed.
A \textbf{cycle} $(a_1\ a_2\ \cdots\ a_k)$ sends each displayed entry
to the next and the last back to the first, fixing all others.

\subsection{Reading and composing bijections}
A two-row display places each input above its output:
\[
\sigma=\begin{pmatrix}1&2&3&4&5\\3&1&2&5&4\end{pmatrix}.
\]
Its output list $(3,1,2,5,4)$ is shorter, but a cycle uses a different
convention: $(1\ 3\ 2)$ means $1\mapsto3\mapsto2\mapsto1$.
Thus $\sigma=(1\ 3\ 2)(4\ 5)$. Starting the same cycle elsewhere
does not change it: $(1\ 3\ 2)=(3\ 2\ 1)=(2\ 1\ 3)$.
Reversing it does change the map and gives its inverse $(1\ 2\ 3)$.
An omitted entry is fixed; it is not absent from the domain.
\begin{center}
\begin{tikzpicture}[>=Stealth]
\node (a) at (0,1) {$1$}; \node (b) at (1.3,-.5) {$3$};
\node (c) at (-1.3,-.5) {$2$};
\draw[->] (a)--(b); \draw[->] (b)--(c); \draw[->] (c)--(a);
\node (d) at (4,.8) {$4$}; \node (e) at (4,-.6) {$5$};
\draw[->,bend left=30] (d) to (e); \draw[->,bend left=30] (e) to (d);
\end{tikzpicture}
\end{center}
\begin{example}[Why order matters]
Let $\alpha=(1\ 2\ 3)$ and $\beta=(1\ 2)$ in $S_3$.
Follow each input through two arrows:
\[
\begin{array}{c|ccc|cc}
j&\beta(j)&\alpha(\beta(j))& &\alpha(j)&\beta(\alpha(j))\\\hline
1&2&3&&2&1\\2&1&2&&3&3\\3&3&1&&1&2
\end{array}
\]
Thus $\alpha\beta=(1\ 3)$, but $\beta\alpha=(2\ 3)$.
By contrast $(1\ 3\ 2)$ and $(4\ 5)$ act on disjoint sets of entries.
For an entry in the first set, the second cycle does nothing, and conversely;
every other entry is fixed by both. Hence disjoint cycles commute.
\end{example}
\subsection{Following a permutation until it closes}
Why should cycles describe every bijection? Start at one entry and keep
following its image. There are finitely many entries, so eventually one
repeats. Injectivity prevents the path from joining itself anywhere except
its starting point: if $\sigma^r(a)=\sigma^s(a)$ with $0<r<s$, applying
the inverse $r$ times gives $a=\sigma^{s-r}(a)$, an earlier return.
After closing one cycle, begin at an unused entry.
\begin{proposition}[Cycle decomposition]
Every permutation decomposes into disjoint cycles, uniquely up to the
order of the cycles and cyclic rotation of their entries. Every permutation
is a product of transpositions, where a transposition is a two-entry cycle.
\end{proposition}
\begin{proof}
Following successive images of an entry closes a cycle: finiteness forces
a repetition, and injectivity makes the first repetition the starting entry.
Starting again outside the entries already used gives disjoint cycles and
exhausts the finite set. Each cycle is the orbit of any one of its entries,
with its order prescribed by the permutation, proving the stated uniqueness.
Finally
\[
(a_1\ a_2\ \cdots\ a_k)
=(a_1\ a_k)(a_1\ a_{k-1})\cdots(a_1\ a_2).
\]
Checking the images of the displayed entries proves this identity; applying
it to each disjoint cycle gives a transposition decomposition.
\end{proof}
For example the permutation with list $(2,3,1,5,4)$ is $(1\ 2\ 3)(4\ 5)$,
and equals $(1\ 3)(1\ 2)(4\ 5)$.

\begin{example}[A full decomposition and its inverse]
The list $(4,3,2,5,1,6)$ sends $1\to4\to5\to1$,
$2\to3\to2$, and fixes $6$. Hence
\[
\sigma=(1\ 4\ 5)(2\ 3)=(1\ 5)(1\ 4)(2\ 3),\qquad
\sigma^{-1}=(1\ 5\ 4)(2\ 3).
\]
For a cycle $(a_1\cdots a_k)$, the transposition formula can be verified
without multiplying symbols formally. The rightmost swap takes $a_1$ to
$a_2$, which later swaps leave alone. For $2\leq j<k$, the swap
$(a_1\ a_j)$ takes $a_j$ to $a_1$, and the next takes it to $a_{j+1}$.
The final entry $a_k$ goes to $a_1$. Every unlisted entry stays fixed.
\end{example}
\subsection{Adjacent swaps and the parity question}
An \textbf{adjacent transposition} is $(i\ i+1)$. Swapping positions
$i<j$ can be accomplished by moving the entry at $i$ rightward to $j$
and moving the displaced entry leftward back to $i$. Intermediate entries
return to their original places. This uses $2(j-i)-1$ adjacent swaps:
\[
(i\ j)=(i\ i+1)\cdots(j-2\ j-1)(j-1\ j)
(j-2\ j-1)\cdots(i\ i+1).
\]
For example $(1\ 3)=(1\ 2)(2\ 3)(1\ 2)$. The same permutation
therefore has a one-swap and a three-swap expression. Inserting
$(1\ 2)(1\ 2)$ produces still more expressions. Length is not intrinsic.
We must determine whether its parity is intrinsic: could the same
permutation have both a two-swap and a three-swap expression?

We need parity to be independent of the chosen transpositions. The next
proof uses formal polynomials with integer coefficients only as finite
expressions in indeterminates; it does not assume any polynomial theory.
\begin{theorem}[Parity and sign]
\label{thm:appendix-parity}
Any two transposition expressions for a permutation have the same parity
of length. Define $\operatorname{sgn}(\sigma)=(-1)^r$ using that parity.
Then sign is multiplicative under composition.
\end{theorem}
\begin{proof}
We seek a nonzero expression whose value reacts predictably when two
positions are exchanged. The Vandermonde expression has exactly this
property: every transposition negates it. It can therefore determine
whether a transposition decomposition has even or odd length. Consider
$\Delta(t_1,\ldots,t_n)=\prod_{i<j}(t_j-t_i)$ over the integers.
Swapping adjacent variables $t_i,t_{i+1}$ negates their mutual factor.
For a third index $k<i$, the pair of factors is
$(t_i-t_k)(t_{i+1}-t_k)$; for $k>i+1$ it is
$(t_k-t_i)(t_k-t_{i+1})$. Each pair merely exchanges its factors.
Thus only the mutual factor changes sign and the whole expression is negated. Swapping positions $i<j$ can be done in
$2(j-i)-1$ adjacent swaps, hence also negates $\Delta$.
Thus a decomposition of length $r$ changes $\Delta$ to $(-1)^r\Delta$.
Two decompositions of the same permutation perform the same substitution,
so yield equal expressions. They cannot yield opposite signs: evaluating
at the explicit integer tuple gives
\[
\Delta(1,\ldots,n)=\prod_{i<j}(j-i)>0
\]
(the empty product is $1$), so $\Delta$ is nonzero and not its negative.
Parity is therefore well defined. Concatenating decompositions proves
$\operatorname{sgn}(\sigma\tau)=\operatorname{sgn}(\sigma)
\operatorname{sgn}(\tau)$.
\end{proof}
Call permutations with sign $1$ \textbf{even} and sign $-1$ \textbf{odd}.
A $k$-cycle has sign $(-1)^{k-1}$.
The parity is defined over the integers even when coefficients later lie
in characteristic two; interpreting its signs in that field sends both
$+1$ and $-1$ to $1$ without changing the combinatorial parity theorem.

The identity uses zero swaps and has sign $1$. Reversing a swap expression
gives an expression of the same length for the inverse, since each swap
undoes itself. Therefore $\operatorname{sgn}(\sigma^{-1})
=\operatorname{sgn}(\sigma)$. A $k$-cycle uses $k-1$ swaps; a
permutation with disjoint cycles of lengths $k_1,\ldots,k_r$ has sign
$(-1)^{\sum_j(k_j-1)}$. Fixed points contribute zero to this exponent.
\subsection{Counting inversions instead of choosing swaps}

The Vandermonde argument explains why parity is intrinsic; inversion
counting now supplies an efficient way to compute it.
Parity is known to be independent of a swap decomposition. To compute
it without finding such a decomposition, count the pairs that appear in
the wrong order. Each adjacent correction removes exactly one inversion.

\begin{proposition}[Inversions compute the same sign]
An inversion of $\sigma$ is a pair $i<j$ with $\sigma(i)>\sigma(j)$.
Then $\operatorname{sgn}(\sigma)=(-1)^{\operatorname{inv}(\sigma)}$.
\end{proposition}
\begin{proof}
If the list is not increasing, it has an adjacent inverted pair.
Exchanging that pair decreases the inversion count by one: their mutual
inversion disappears and comparisons with all other entries have the same
total count. After exactly $\operatorname{inv}(\sigma)$ exchanges the list
is increasing, hence is the identity. Parity now gives the formula.
\end{proof}
\begin{example}[Inversions in a complete list]
For $(3,1,2,5,4)$ the inverted position pairs are $(1,2),(1,3),(4,5)$.
There are three, so the permutation is odd. Sorting by adjacent exchanges
makes the count visible:
\[
31254\longrightarrow13254\longrightarrow12354\longrightarrow12345.
\]
Its cycle expression $(1\ 3\ 2)(4\ 5)$ also gives sign
$(-1)^2(-1)=-1$. The two calculations agree for a proved reason, not
because a preferred decomposition was selected.
\end{example}
In the inversion proof, a third entry to the left compares with both
exchanged entries before and after; the number of these two inversions is
unchanged. A third entry to the right has the same property. Only the
comparison between the adjacent entries reverses. Moreover an unsorted
list must have adjacent entries out of order: if every neighbor increased,
transitivity would make the entire list increasing.
\subsection{Permutation matrices connect the notation to linear maps}
The \textbf{permutation matrix} $P_\sigma$ has column $j$ equal to
$e_{\sigma(j)}$. Applying it to basis vectors shows
$P_\sigma P_\tau=P_{\sigma\tau}$ and $P_\sigma^{-1}=P_{\sigma^{-1}}$.

\begin{example}[Which way does a permutation matrix move coordinates?]
For $\sigma=(1\ 2\ 3)$ the columns are $e_2,e_3,e_1$, hence
\[
P_\sigma=\begin{pmatrix}0&0&1\\1&0&0\\0&1&0\end{pmatrix},\qquad
P_\sigma\begin{pmatrix}x_1\\x_2\\x_3\end{pmatrix}
=\begin{pmatrix}x_3\\x_1\\x_2\end{pmatrix}.
\]
The old coordinate $x_j$ moves to position $\sigma(j)$; the new
coordinate in position $i$ is $x_{\sigma^{-1}(i)}$.
In the Leibniz sum only the selection $(2,3,1)$ has a nonzero product.
Its sign is $+1$, so $\det P_\sigma=1$. This example also explains
why a permutation and its inverse can appear in different coordinate
conventions without changing the determinant sign.
\end{example}
\section{Constructing the Determinant Form}

We have constructed a sign for each permutation and proved how signs
behave under composition. This is the missing coefficient in the
determinant formula: multilinearity chooses one basis vector in each slot,
alternation removes repeated choices, and permutation sign records the
orientation of the surviving orderings.

Let $(v_1,\ldots,v_n)$ be a basis of $V$ and write
$x_j=\sum_i a_{ij}v_i$. The slot-by-slot uniqueness argument of
\cref{thm:determinant-form} uses only multilinearity and alternation,
so works over $F$: each alternating form is determined by its value on
this basis. We now give the field-valid existence argument in full.
\begin{theorem}[Existence over an arbitrary field]
The formula
\[
\omega(x_1,\ldots,x_n)=\sum_{\sigma\in S_n}
\operatorname{sgn}(\sigma)\prod_{j=1}^n a_{\sigma(j),j}
\]
defines the unique alternating $n$-linear form with
$\omega(v_1,\ldots,v_n)=1$.
\end{theorem}
\begin{proof}
Each product uses one coordinate from each input, so
distributing sums or scalar multiples in one column proves multilinearity.
If $x_r=x_s$ for $r\ne s$, pair $\sigma$ with $\sigma\circ(r\ s)$.
This is a fixed-point-free pairing. The products agree since the two
coordinate columns agree, and the signs are opposite by parity. Their
sum is zero in every field, including characteristic two where it is
twice the same element. On the given basis, all terms except the identity
term vanish, and that term is $1$. For uniqueness, expanding an arbitrary
form gives the same sum: repeated tuples vanish, and the surviving tuples
are permutations with their forced signs.
\end{proof}
In the standard basis this is $\det A$. It agrees with the main-text
determinant when $F=\R$ by uniqueness, and gives
$\det P_\sigma=\operatorname{sgn}(\sigma)$ interpreted in $F$.
The uniqueness proof of multiplicativity, the invertibility criterion,
the transpose formula, and basis independence in Chapter~10 remain valid
over $F$. No order or positivity was used in them.

\begin{proposition}[Laplace expansion from the permutation sum]
The minor and cofactor definitions and both expansions in
\cref{thm:laplace-expansion} are valid over $F$.
\end{proposition}
\begin{proof}
Group the Leibniz terms by the row $i=\sigma(j)$ selected in the fixed
column $j$. Delete row $i$ and column $j$, preserving the order of all
remaining rows and columns. The residual choices are precisely the
permutations of this minor. Moving row $i$ and column $j$ to the first
positions uses $(i-1)+(j-1)$ adjacent swaps; their sign is $(-1)^{i+j}$.
The contribution of this group is therefore $a_{ij}(-1)^{i+j}M_{ij}$.
Summing over $i$ gives column expansion, and transpose invariance gives
row expansion. This is the permutation counterpart of the alternating-form
proof of \cref{thm:laplace-expansion}, valid over every field.
\end{proof}
For example the sparse matrix of Chapter~10 has determinant
$-6(8-3)=-30$, interpreted in $F$. The same calculation remains valid
if this field element becomes zero: $-30=0$ in characteristics $2,3,5$,
whereas $-30=5$ in characteristic $7$. The matrix is singular in the
first three cases and invertible in the last. Interpret the scalar in
$F$ before using the determinant criterion.

\paragraph{A normalized form versus an endomorphism determinant.}
Let $\omega_e$ be normalized to $1$ on a basis $e$, and let the columns
of $P$ express a new basis $e'$ in $e$ coordinates. Then
$\omega_e(e'_1,\ldots,e'_n)=\det P$, so uniqueness of the normalized form gives
\[
\omega_{e'}=(\det P)^{-1}\omega_e.
\]
For a fixed endomorphism $T:V\to V$, its matrix changes instead from
$A$ to $P^{-1}AP$; multiplicativity makes their determinants equal.
For a map between distinct spaces there are separate source and target
basis changes, so there is no independent scalar determinant without
normalizations. The top exterior-power action in
\cref{thm:top-exterior-determinant} will give a fourth equivalent viewpoint
alongside normalized alternation, Leibniz, and recursive Laplace expansion.

\Needspace{10\baselineskip}
\section*{Exercises}
\addcontentsline{toc}{section}{Exercises}

\begin{exercise}
Find disjoint cycles, a transposition decomposition, sign, inversion count,
and determinant of the permutation matrix for the list $(2,5,4,1,3)$.
\end{exercise}
\begin{hexercise}{app-c-2}
Define the adjugate by $(\operatorname{adj}A)_{ji}=C_{ij}(A)$.
Using cofactor expansion of a matrix with one replaced row, prove
$A\operatorname{adj}A=(\det A)I$. Derive the inverse formula when
$\det A\ne0$. For $Ax=b$, use column multilinearity to derive
$x_j=\det(A_1,\ldots,b,\ldots,A_n)/\det A$ (Cramer's rule).
\end{hexercise}

\begin{exercise}
In $S_5$, compute both $(1\ 4\ 2)(2\ 3\ 5)$ and
$(2\ 3\ 5)(1\ 4\ 2)$ by making a table of all five images. Give
inverse lists and cycle expressions. Which entries witness noncommutation?
\end{exercise}
\begin{exercise}
Give three distinct transposition expressions for $(1\ 4)$ and count
their lengths. Explain what parity says, and what it does not say, about
the smallest possible length of an expression.
\end{exercise}
\begin{hexercise}{app-c-5}
For the reversing list $(n,n-1,\ldots,1)$ find the inversion count and
sign. Compute its permutation matrix determinant for $n=4,5,6$.
\end{hexercise}
\begin{hexercise}{app-c-6}
Over $\mathbb F_2$, compute the determinant of
$\left(\begin{matrix}1&1\\1&0\end{matrix}\right)$.
Explain why the answer gives an invertibility test but no positive or
negative orientation. Distinguish permutation parity from its field-valued sign.
\end{hexercise}
